\section{Variable-Granularity Object Storage}
\label{sec:layout}

\subsection{An independent choice of object granularity}

Object granularity controls how file data becomes uploadable objects and
how much work a subsequent access requires. It is useful in an ordinary
flat volume as well as in a workspace or an immutable snapshot; neither
delta-based workspaces nor Packed Metadata is a prerequisite. A deployment
can retain the fixed client block size and the backend's default layout.
Selecting another granularity changes this data path without changing the
namespace's inheritance or deletion rules.

There are several distinct boundaries. Let $g$ denote the client block size,
and $L_o$ the finalized upload length after applicable
compression, headers, and framing. The client then chooses which object
ranges to fetch. On a compatible backend, $B$ denotes the erasure-coding
block size and $q$ the physical verification/read quantum within a shard.
These quantities describe different work: object formation, transferred
bytes, encoding, and physical verification. Changing one does not
necessarily change the others.

\subsection{Forming objects from file updates}

The ordinary data path locates a file offset within a logical chunk and
records written content in slices. A chunk is an addressing range, rather
than a remote object. The uploader partitions each slice's data into blocks
using $g$ and stores them under the slice identity and block index. Thus a
10 MiB slice with $g=4\MiB$ can produce raw blocks of 4, 4, and 2 MiB.
Compression and framing subsequently determine each $L_o$; a short tail
need not occupy a full block. Flush boundaries and partial writes also
affect the objects actually formed.

The reader resolves the visible slice intervals and translates them into
object offsets using the same geometry. A new slice can cover an old
interval while untouched intervals still refer to older objects. Such an
update need not rewrite the complete old object. Compaction, upload retries,
and the accumulation of short slices can nevertheless add work beyond the
newly written bytes.

The present ordinary path defaults to 64 MiB chunks and $g=4\MiB$.
Its client block size is a fixed configuration, with a chunk size that is
a nonzero multiple of $g$. The proposed 1/4/8/16 MiB choices express a
design space for object formation; an online controller that selects among
them from access behavior is not implemented in this path. Readers currently
reconstruct block addresses from the configured geometry, so changing $g$
for existing data requires a compatible representation or migration.

\subsection{Trading request count against transferred work}

Larger objects can reduce the number of requests needed to transfer a long
contiguous slice. Smaller objects can reduce excess fetching when decoding,
authentication, or cache policy requires a complete object. Table~\ref{tab:object-granularity}
illustrates this tradeoff for one uncompressed, aligned 64 MiB slice and a
read of file interval $[10,12)\MiB$. These are arithmetic examples, not
measurements or a policy recommendation.

\begin{table}[t]
\centering
\begin{tabular}{rrr}
\toprule
$g$ & Object count & \shortstack{Whole-object\\bytes fetched}\\
\midrule
1 MiB & 64 & 2 MiB\\
4 MiB & 16 & 4 MiB\\
8 MiB & 8 & 8 MiB\\
16 MiB & 4 & 16 MiB\\
\bottomrule
\end{tabular}
\caption{Idealized object formation and a 2 MiB partial read. The last
column assumes complete fetches of the intersecting objects, with no cache
hits or compression; it does not describe every range-read path.}
\label{tab:object-granularity}
\end{table}

Range GET can retrieve only part of an object when its format permits
independent access. The ordinary store supports page-aligned range reads
for eligible small requests, while compressed blocks or a whole-object
fetch policy can require larger transfers. A packed frame supplies another
decoding and integrity boundary inside a snapshot object. Its size is
neither $g$ nor the backend's $q$. Caching and concurrent-read coalescing
further change how many remote requests reach the backend.

We therefore separate network read amplification, upload amplification,
and physical read amplification. Their numerators are,
respectively, object response bytes, uploaded bytes, and storage-device
bytes; their denominators are the corresponding application bytes served
or written. Request counts, compression work, and cache occupancy remain
separate costs. Smaller objects do not guarantee lower write amplification,
and a large object does not imply a complete download for every small read.

\subsection{Optional backend layout}

A compatible RustFS research backend provides a further choice after the
object has formed. The client may send an encoding-block-size hint
$h(L_o)$ based on the finalized upload length. This interface leaves $g$
and multipart part sizes unchanged. Omitting the hint retains the backend's
default selection. The backend validates an enabled request and persists
the effective $B$; an emitted header alone does not establish which layout
was used. Multipart initiation uses the complete object's length, and
retries retain that object's decision.

The current hint modes are disabled, fixed, and automatic. Fixed candidates
are 64 KiB, 256 KiB, 1 MiB, and 4 MiB. The automatic hint policy omits the
hint for $L_o\leq1\MiB$ and requests $B=1\MiB$ otherwise. This is a
threshold rule rather than access-adaptive object formation. Because the
inspected backend's default $B$ is also 1 MiB, this rule may leave the
effective encoding block size unchanged.

For $d$ data shards and $p$ parity shards, let $s$ be one complete stripe's
logical shard size. A separately enabled backend policy can select a
physical verification/read quantum subject to
\begin{equation}
q\mid s, \qquad qd\mid B.
\end{equation}
These constraints preserve stripe alignment while permitting smaller
verification units. Both hint admission and new $q$ framing have their own
compatibility gates. Reads and healing use persisted object geometry even
after write settings change.

For a nonempty shard interval starting at $a$ with length $\ell$, the number
of verification/read quanta intersected is
\begin{equation}
n_q(a,\ell)=
\left\lceil\frac{a+\ell}{q}\right\rceil-
\left\lfloor\frac{a}{q}\right\rfloor.
\end{equation}
Ignoring tails and batching, the covered data is approximately $n_q q$
plus checksums, summed over shards actually read. The application request
cannot be substituted directly for $\ell$: the client range, compression,
and coding determine that interval. Smaller $q$ can reduce over-read while
increasing checksum overhead. Fixed $(d,p)$ still has nominal redundancy
$(d+p)/d$; a large $B$ does not imply a fully padded tail stripe.

\subsection{Composition and scope}

A workspace changes which objects a view references; Packed Metadata
changes how immutable records and payload locations are obtained. Either
can use default backend layouts. Conversely, a flat volume can configure
$g$ and eligible backend hints without either mechanism. Snapshot producers
and workspace upper writers retain their own object formats and upload
contracts, so one path's block geometry cannot be assumed for all paths.
Section~\ref{sec:implementation} identifies the current integration limits.
It also records the ordinary path's coupling between enabled hints and
whole-object caching, which must be held constant when isolating $B$ or $q$.
